\documentclass[11pt]{article}
\usepackage[utf8]{inputenc}
\usepackage{amsmath,amssymb}
\usepackage{hyperref}
\title{Scalable Neurodegeneration Protein Aggregation for Large-Scale Settings}
\author{Anonymous}
\date{}
\begin{document}
\maketitle

\begin{abstract}
This paper studies Neurodegeneration Protein Aggregation. Grounded in a real, citable literature \cite{ref1,ref2,ref3,ref4,ref5,ref6,ref7,ref8},
we propose a focused method and report \textbf{illustrative} experiments.
All references below are real and verifiable (OpenAlex/arXiv); the reported
numbers are simulated to illustrate intended behaviour.
\end{abstract}

\section{Introduction}
Neurodegeneration Protein Aggregation is an active research area. This work targets a concrete gap identified from
the recent literature and contributes a method together with an evaluation protocol.

\section{Related Work (real literature)}
The following works, retrieved from public bibliographic APIs, frame our study:
\begin{itemize}
\item Haass et al. (2007) — "Soluble protein oligomers in neurodegeneration: lessons from the Alzheimer's amyloid β-peptide", Nature Reviews Molecular Cell Biology (cited 4731).
\item Wang et al. (2015) — "Tau in physiology and pathology", Nature reviews. Neuroscience (cited 2240).
\item Jarrett et al. (1993) — "Seeding “one-dimensional crystallization” of amyloid: A pathogenic mechanism in Alzheimer's disease and scrapie?", Cell (cited 2177).
\item Kopito (2000) — "Aggresomes, inclusion bodies and protein aggregation", Trends in Cell Biology (cited 1922).
\item Ross et al. (2010) — "Huntington's disease: from molecular pathogenesis to clinical treatment", The Lancet Neurology (cited 1740).
\item Lashuel et al. (2012) — "The many faces of α-synuclein: from structure and toxicity to therapeutic target", Nature reviews. Neuroscience (cited 1677).
\end{itemize}

\section{Method}
We decompose the problem across shards with a communication-efficient aggregation rule that keeps overhead sublinear in scale. We instantiate this idea for Neurodegeneration Protein Aggregation and analyse its cost/quality trade-off.

\section{Experiments (Illustrative)}
\begin{center}
\begin{tabular}{lcc}
System & Primary Metric & Efficiency \\ \hline
Strong baseline & 0.812 & 1.00$\times$ \\
Ablation & 0.828 & 0.92$\times$ \\
\textbf{Ours} & \textbf{0.851} & \textbf{0.71$\times$}
\end{tabular}
\end{center}
\emph{Metrics simulated to illustrate the intended effect; not measured.}

\section{Conclusion}
We connected a real literature to a concrete method for Neurodegeneration Protein Aggregation. Future work will
validate the illustrative results with real experiments.

\begin{thebibliography}{9}
\bibitem{ref1} Christian Haass, Dennis J. Selkoe. Soluble protein oligomers in neurodegeneration: lessons from the Alzheimer's amyloid β-peptide. \emph{Nature Reviews Molecular Cell Biology}, 2007. DOI:10.1038/nrm2101
\bibitem{ref2} Yipeng Wang, Eckhard Mandelkow�. Tau in physiology and pathology. \emph{Nature reviews. Neuroscience}, 2015. DOI:10.1038/nrn.2015.1
\bibitem{ref3} Joseph T. Jarrett, Peter T. Lansbury. Seeding “one-dimensional crystallization” of amyloid: A pathogenic mechanism in Alzheimer's disease and scrapie?. \emph{Cell}, 1993. DOI:10.1016/0092-8674(93)90635-4
\bibitem{ref4} Ron R. Kopito. Aggresomes, inclusion bodies and protein aggregation. \emph{Trends in Cell Biology}, 2000. DOI:10.1016/s0962-8924(00)01852-3
\bibitem{ref5} Christopher A. Ross, Sarah J. Tabrizi. Huntington's disease: from molecular pathogenesis to clinical treatment. \emph{The Lancet Neurology}, 2010. DOI:10.1016/s1474-4422(10)70245-3
\bibitem{ref6} Hilal A. Lashuel, Cassia Overk, Abid Oueslati et al.. The many faces of α-synuclein: from structure and toxicity to therapeutic target. \emph{Nature reviews. Neuroscience}, 2012. DOI:10.1038/nrn3406
\bibitem{ref7} Mathias Jucker, Lary C. Walker. Self-propagation of pathogenic protein aggregates in neurodegenerative diseases. \emph{Nature}, 2013. DOI:10.1038/nature12481
\bibitem{ref8} Ted M. Dawson, Valina L. Dawson. Molecular Pathways of Neurodegeneration in Parkinson's Disease. \emph{Science}, 2003. DOI:10.1126/science.1087753
\end{thebibliography}
\end{document}
